\documentclass[11pt]{article}
\usepackage[margin=1in]{geometry}
\usepackage{amsmath,amssymb,amsthm,mathtools}
\usepackage[hidelinks]{hyperref}
\usepackage{booktabs,array,longtable,pdflscape}

\newtheorem{theorem}{Theorem}[section]
\newtheorem{lemma}[theorem]{Lemma}
\newtheorem{proposition}[theorem]{Proposition}
\newtheorem{corollary}[theorem]{Corollary}
\theoremstyle{remark}
\newtheorem{remark}[theorem]{Remark}
\newcommand{\splus}{s^+}
\newcommand{\sminus}{s^-}
\newcommand{\one}{\mathbf 1}
\newcommand{\kap}{\kappa}
\newcommand{\sig}{\sigma}

\title{Positive Square Energy of Every Tetracyclic Graph}
\author{AI-generated manuscript; no human author is claimed}
\date{31 July 2026}

\begin{document}
\maketitle

\begin{abstract}
Let $G$ be a finite simple connected graph with $n$ vertices and $n+3$
edges.  We prove that the sum $\splus(G)$ of the squares of its positive
adjacency eigenvalues is at least $n$.  The positive cyclomatic ranks of the
cyclic blocks partition four as $1+1+1+1$, $2+1+1$, $2+2$, $3+1$, or $4$.
The first four alternatives are closed by a sharp block-additive doubly
nonnegative bound and explicit induced-territory packets for its residuals.
In the last alternative, suppression of degree-two paths gives exactly
seventeen loopless $2$-connected multigraph kernels.  We include the DNN
correlation dual, exact path elimination, the kernel classification, and
complete count ledgers for the exact finite certificates covering all
seventeen kernels.  Parity monotonicity turns each finite frontier into an
all-length proof, and one-vertex additivity includes arbitrary rooted trees.
Some subsidiary estimates are strict, but boundary certificates occur, so no
strict global assertion is made.
\end{abstract}

\section{Statement and spectral reductions}

All graphs are finite, simple, and undirected.  Let
$\lambda_1,\ldots,\lambda_n$ be the adjacency eigenvalues and put
\[
 \splus(G)=\sum_{\lambda_i>0}\lambda_i^2,
 \qquad \sminus(G)=\sum_{\lambda_i<0}\lambda_i^2,
 \qquad D(G)=\splus(G)-\sminus(G).
\]
Since $\operatorname{tr}A(G)^2=2|E(G)|$, a tetracyclic graph satisfies
\begin{equation}\label{eq:trace}
 \splus(G)+\sminus(G)=2n+6,
 \qquad \splus(G)=n+3+\frac12D(G).
\end{equation}

\begin{theorem}\label{thm:main}
Every finite simple connected graph $G$ with $|V(G)|=n$ and
$|E(G)|=n+3$ satisfies
\[
 \boxed{\splus(G)\geq n}.
\]
\end{theorem}

We use square-energy credit $\sig(H)=\splus(H)-|V(H)|$.  The following
standard pinching inequality permits structural repairs.

\begin{lemma}[Induced superadditivity]\label{lem:super}
If $V(G)=V_1\mathbin{\dot\cup}\cdots\mathbin{\dot\cup}V_r$, then
\[
 \splus(G)\geq\sum_{j=1}^r\splus(G[V_j]),
 \qquad
 \sig(G)\geq\sum_{j=1}^r\sig(G[V_j]).
\]
\end{lemma}

\begin{proof}
For a real symmetric matrix $A$,
\[
 \operatorname{tr}(A_+^2)=
 \max_{X\succeq0}\bigl(2\operatorname{tr}(AX)-\operatorname{tr}(X^2)\bigr).
\]
Insert the block diagonal matrix whose blocks are the positive parts of the
principal matrices $A[V_j]$.  Cross edges make no contribution to the trace.
\end{proof}

A nonempty connected tree $T$ is bipartite, so
$\splus(T)=|E(T)|=|V(T)|-1$ and $\sig(T)=-1$.  We shall use the following
attachment-uniform credits, all proved in the cited lower-rank manuscripts.

\begin{lemma}[Lower-rank packets]\label{lem:packets}
Arbitrary rooted trees may be attached at arbitrary vertices in every item.
\begin{enumerate}
\item Every connected bicyclic graph has positive credit; a bipartite one has
credit exactly one.  A two-triangle cactus and $\Theta(1,2,2)$ have credit
strictly greater than one.
\item Every connected tricyclic graph has nonnegative credit.
\item A triangular unicyclic cactus has positive credit.  A pentagonal
unicyclic cactus has credit at least $2-\sqrt5$.
\item A triangle--pentagon bicyclic cactus has credit greater than
$3-\sqrt5$, and a three-cycle packet $C_3,C_3,C_5$ whose triangles share a
cut has credit greater than $4-\sqrt5$.
\end{enumerate}
\end{lemma}

\begin{proof}
The bicyclic and tricyclic statements are the complete lower-rank theorems
\cite{BicyclicGraphs,TricyclicGraphs}.  The sharper cactus credits follow from
the signed Coulson formula after rooted-tree matching-message elimination
\cite{TetracyclicCactus}.  They are quoted with exactly the margins used
below; no edge-addition monotonicity is involved.
\end{proof}

\section{The DNN dual and exact paths}

For an edge-containing graph $H$, define
\begin{equation}\label{eq:kappa-primal}
 \kap(H)=\sup_{\substack{M\succeq0,\ M\geq0\\\one^TM\one>0}}
 \frac{4\bigl(\sum_{uv\in E(H)}\sqrt{M_{uv}}\bigr)^2}
 {\one^TM\one}.
\end{equation}

\begin{lemma}[DNN inequality and correlation dual]\label{lem:dnn}
For every graph $H$,
\begin{equation}\label{eq:dnn-dual}
 \sminus(H)\leq\kap(H),
 \qquad
 \kap(H)=\min_{\substack{R\succeq0\\R_{vv}=1}}
 \sum_{uv\in E(H)}\frac{2}{1-R_{uv}},
\end{equation}
with boundary values interpreted by limits.  If $H=H_1\vee H_2$ is a
one-vertex sum, then
\[
 \kap(H)=\kap(H_1)+\kap(H_2).
\]
In particular, $\kap(T)=|E(T)|$ for every tree.
\end{lemma}

\begin{proof}
Write $A=P-B$ for the positive and negative spectral parts.  The Schur product
theorem makes $M=B\circ B$ positive semidefinite and entrywise nonnegative.
Moreover
\[
 \one^TM\one=\operatorname{tr}B^2=\sminus(H),
 \qquad
 \frac{\sminus(H)}2=-\sum_{uv\in E(H)}B_{uv}
 \leq\sum_{uv\in E(H)}\sqrt{M_{uv}}.
\]
Substitution in \eqref{eq:kappa-primal} proves $\sminus\leq\kap$.

For weak duality, if $R$ is a correlation matrix, then
\[
 \one^TM\one-
 \sum_{uv\in E(H)}2(1-R_{uv})M_{uv}
 =\langle R,M\rangle+
 2\sum_{\substack{u<v\\uv\notin E(H)}}(1-R_{uv})M_{uv}\geq0.
\]
Weighted Cauchy--Schwarz gives the upper bound in \eqref{eq:dnn-dual}.
For the reverse bound normalize $\langle J,M\rangle=1$, use
\[
 (\sum_e\sqrt{x_e})^2
 =\min_{a_e>0,\ \sum a_e=1}\sum_e\frac{x_e}{a_e},
\]
and apply conic minimax with
$(\mathcal S_+\cap\mathcal N)^*=\mathcal S_++\mathcal N$; an
$\varepsilon$ truncation treats boundary points.

For a one-vertex sum, restriction of a feasible correlation matrix gives one
inequality.  Conversely, realize the two correlation matrices by unit Gram
vectors with their common cut vector equal and place their orthogonal
components in orthogonal subspaces.  This glues feasible matrices and proves
the opposite inequality.  Iterating $\kap(K_2)=1$ proves the tree formula.
\end{proof}

\begin{lemma}[Exact path elimination]\label{lem:path}
Replace a kernel edge $uv$ by a path of length $\ell$.  For fixed endpoint
correlation $r=R_{uv}$, the exact excess of its contribution over $\ell$ is
\begin{equation}\label{eq:path-cost}
 f_\ell(r)=\ell\tan^2\!\left(
 \frac{\arccos((-1)^\ell r)}{2\ell}\right).
\end{equation}
For fixed $r$ and fixed parity, $f_{\ell+2}(r)\leq f_\ell(r)$, strictly unless
the transformed endpoint angle is zero.  If a subdivided kernel has total path
length $L$, then
\begin{equation}\label{eq:kernel-min}
 \kap(H)-L=\min_R\sum_{P}f_{|P|}(R_{u_Pv_P}),
\end{equation}
where $R$ ranges over correlation matrices on the branch vertices.
\end{lemma}

\begin{proof}
Represent $R$ by unit vectors and alternately negate the vectors along the
path.  The transformed endpoint angle is
$\beta=\arccos((-1)^\ell r)$.  The spherical triangle inequality makes the
sum of successive angles at least $\beta$, and convexity of
$\sec^2(x/2)$ makes equal angles optimal.  Equally spaced planar vectors on a
shortest arc attain equality, yielding
$\ell\sec^2(\beta/(2\ell))=\ell+f_\ell(r)$.  The constructions for distinct
paths share only their fixed endpoint vectors and are therefore compatible.
Finally, for $h_s(x)=s\tan^2(x/s)$,
\[
 \frac{d}{ds}h_s(x)=
 \tan z\sec^2z\bigl(\sin z\cos z-2z\bigr)<0,
 \qquad z=x/s,
\]
because $0<\sin z\cos z<z<2z$.  This proves fixed-parity monotonicity.
\end{proof}

If a connected rank-$r$ core has $L$ edges, it has $L-r+1$ vertices.
Consequently, for rank four,
\begin{equation}\label{eq:budget}
 \kap(H)\leq L+3
 \quad\Longrightarrow\quad
 \splus(H)\geq2L-(L+3)=L-3=|V(H)|.
\end{equation}
Tree attachments preserve this calculation, since each new tree edge adds one
to $L$, $|V|$, and $\kap$.

\section{Block partitions and all multiblock residuals}

A block is a maximal connected subgraph without a cut vertex; bridges count
as two-vertex blocks.  Cyclomatic rank is additive over blocks.  Suppression
and all degree statements below are internal to the relevant block.

\begin{proposition}[Block alternatives]\label{prop:block-alternatives}
The positive ranks of the cyclic blocks of a connected tetracyclic graph are
exactly
\[
 1+1+1+1,\qquad2+1+1,\qquad2+2,\qquad3+1,\qquad4.
\]
Rank-one blocks are cycles, and rank-two blocks are subdivisions of three
parallel kernel edges, i.e. theta graphs.
\end{proposition}

\begin{proof}
The positive block ranks are positive integers summing to four.  A rank-one
$2$-connected block is a cycle.  After suppressing degree-two vertices of a
rank-two block, the identity $\sum_v(d(v)-2)=2$ and $2$-connectivity force two
degree-three vertices joined by three parallel edges.  Suppression creates no
loop: a suppressed loop would be a cycle attached to the rest at one branch
vertex, making that vertex a cut.
\end{proof}

For a theta with path lengths $a,b,c$, define its exact DNN excess
\begin{equation}\label{eq:theta-excess}
 \Delta(a,b,c)=\min_{0\leq\theta\leq\pi}
 \left(
 \sum_{\ell\ \mathrm{even}}\ell\tan^2\frac{\theta}{2\ell}
 +\sum_{\ell\ \mathrm{odd}}\ell\tan^2\frac{\pi-\theta}{2\ell}
 \right).
\end{equation}
For a cycle let $\epsilon_q=0$ for even $q$ and
$\epsilon_q=q\tan^2(\pi/(2q))$ for odd $q$.  Lemma~\ref{lem:path} gives
$\kap(\Theta)=|E(\Theta)|+\Delta$ and
$\kap(C_q)=q+\epsilon_q$.

The elementary tangent classification needed here is
\begin{equation}\label{eq:theta-class}
 \Delta(a,b,c)>1
 \Longleftrightarrow (a,b,c)=(1,2,r)\ (r\geq2),
 \qquad
 \max\Delta=\Delta(1,2,2)=\frac{\sqrt{17}-1}{2},
\end{equation}
and, away from $(1,2,2)$ in the exceptional family,
$\Delta(1,2,r)<4/3$.  Also
\begin{equation}\label{eq:cycle-values}
 \epsilon_3=1,
 \qquad \epsilon_5=5-2\sqrt5<\frac23,
 \qquad \epsilon_q<\frac25\quad(q\geq7\text{ odd}).
\end{equation}
Here is the full classification argument.  Put $x=\theta/2$ and
$h_s(y)=s\tan^2(y/s)$.  For fixed $0<y<\pi/2$,
\begin{equation}\label{eq:h-monotone}
 \frac{d}{ds}h_s(y)=
 \frac{\tan(y/s)}{\cos^2(y/s)}
 \left(\sin(y/s)\cos(y/s)-\frac{2y}{s}\right)<0.
\end{equation}
Thus an upper-bound test may replace an even path by length two and a nonunit
odd path by length three.  Simplicity permits at most one unit path.

Suppose first that there is no unit path, and let $e$ be the number of even
paths.  It is enough to test
\[
 F_e(x)=2e\tan^2(x/2)+3(3-e)
              \tan^2((\pi/2-x)/3).
\]
For $e=0,3$, use respectively $x=\pi/2,0$.  For the mixed cases use
\[
 F_1(\pi/4)<1,
 \qquad F_2(\pi/6)<1.
\]
The identities $\tan(\pi/8)=\sqrt2-1$ and
$\tan(\pi/12)=2-\sqrt3$, together with
$\tan(\pi/9)<\tan(\pi/8)$, bound the two values by
$48-4\sqrt2-24\sqrt3$ and $37-6\sqrt2-16\sqrt3$; successive squaring of
positive sides proves both are below one.  Hence $\Delta<1$ in this case.

Now suppose one path has length one.  Its contribution is $\cot^2x$.  If the
other paths are both odd, $x=\pi/2$ gives zero.  If exactly one is even and
that even length is at least four, $x=\pi/3$ gives
\[
 \Delta\leq\frac13+4\tan^2\frac\pi{12}
                    +3\tan^2\frac\pi{18}<1;
\]
if both are even and at least four, the same test gives
\[
 \Delta\leq\frac13+8\tan^2\frac\pi{12}<1.
\]
Indeed $\tan(\pi/18)<\tan(\pi/12)<1/3$, and both comparisons reduce to
$48\sqrt3>83$, equivalently $6912>6889$ after squaring.  Therefore the only
remaining sorted triples are $(1,2,r)$, $r\geq2$.  With
$u=\tan(x/2)$, the first two terms obey
\[
 \cot^2x+2\tan^2(x/2)
 =\frac1{4u^2}-\frac12+\frac94u^2\geq1.
\]
The third term is positive in the interior; at either endpoint where it can
vanish, the displayed expression diverges or equals two.  Compactness makes
the minimum strictly greater than one.  This proves the equivalence in
\eqref{eq:theta-class}.

For $r=2$ the complete objective is
\[
 \frac1{4u^2}-\frac12+\frac{17}{4}u^2,
\]
whose minimum, at $u^2=1/\sqrt{17}$, is
$(\sqrt{17}-1)/2$.  For even $r\geq4$, \eqref{eq:h-monotone} makes the
objective pointwise smaller.  For odd $r\geq3$, testing $x=\pi/3$ and using
$\tan(\pi/18)<1/3$ gives $\Delta<4/3$; for even $r\geq4$ the same test gives
$1+4(2-\sqrt3)^2<4/3$.  Thus $(1,2,2)$ is the global maximum and every other
exceptional row has the asserted strict $4/3$ bound.  Finally,
$\epsilon_q=h_q(\pi/2)$ decreases through odd $q$ by
\eqref{eq:h-monotone}; direct evaluation gives $\epsilon_3=1$ and
$\epsilon_5=5-2\sqrt5$, while
$\sin t<t$, $\cos t>1-t^2/2$, and $\pi<22/7$ at $t=\pi/14$ give
$\epsilon_7<2/5$.  This proves \eqref{eq:cycle-values} and reproduces the
complete theta--cycle classification from the tricyclic proof
\cite{TricyclicGraphs}.

\begin{proposition}[Multiblock closure]\label{prop:multiblock}
Every tetracyclic graph having at least two positive-rank cyclic blocks
satisfies $\splus(G)\geq|V(G)|$.
\end{proposition}

\begin{proof}
We give the complete residual ledger.  In each direct row, block additivity of
$\kap$, together with \eqref{eq:budget}, proves the claim.  The last column
names the only rows not direct and the packet that closes them.
\begin{center}
\begin{tabular}{p{1.7cm}p{4.5cm}p{6.4cm}}
\toprule
block ranks&direct DNN criterion&exact residual packet\\
\midrule
$1^4$&$\sum_i\epsilon_{q_i}\leq3$&
$\{3,3,3,q\}$ with odd $q$, and $\{3,3,5,5\}$; the complete cactus
territory theorem closes both families\\
$2+1+1$&$\Delta+\epsilon_p+\epsilon_q\leq3$&
$\Theta(1,2,r)+C_3+C_3$ and
$\Theta(1,2,2)+C_3+C_5$\\
$2+2$&$\Delta_1+\Delta_2\leq3$&
two copies of $\Theta(1,2,2)$\\
$3+1$&$e_B+\epsilon_q\leq3$ for each rank-three certificate&
the canonical doubled-triangle and doubled-$C_4$ class $111$ rows, and the
all-odd $K_4$ rows with at most one long path\\
\bottomrule
\end{tabular}
\end{center}

For $1^4$, the exact cycle calculation leaves precisely the two displayed
families: $\epsilon_3=1$, and among two nontriangular odd cycles only the pair
$5,5$ has excess sum greater than one.  Maximum-packing territories,
shared-cut incidence trees, and the exact twenty-type $C_3,C_3,C_5,C_5$
ledger close them, including arbitrary connectors and rooted branches
\cite{TetracyclicCactus}.

For $2+1+1$, equations \eqref{eq:theta-class}--\eqref{eq:cycle-values} leave
exactly the two rows shown.  In $\Theta(1,2,r)+C_3+C_3$, write the theta paths
as the edge $xy$, the length-two path $xay$, and the $x$--$y$ path $P$ of
length $r$.  Open an internal vertex $v$ of $P$.  Its territory $R$ consists
of $v$, every off-theta component whose block-cut route first meets the theta
at $v$, and every rooted branch based there.  The complementary induced
territory $H$ is connected through the other two theta arms and contains the
intrinsic triangle $xayx$; the remnants of $P-v$ are trees in $H$.  Every cut
vertex belongs to exactly one side, and every component at that cut follows
its unique block-cut-tree route.

Let $k$ be the number of external triangles owned by $R$.  For $k=1$, $R$ is
a triangular packet of credit $>0$ and $H$ is a two-triangle packet of credit
$>1$; for $k=2$ these credits are reversed.  Let $k=0$.  Then $R$ is one
nonempty tree of credit $-1$.  If the three triangles retained in $H$ have
packing number at most two---in particular, if the external triangles share a
cut---their normalized Sachs polynomial has strictly negative imaginary part:
singleton triangle terms are negative imaginary and, since no three cycles
are vertex-disjoint, all other cycle-collection terms are real.  Hence
$D(H)>0$, and because $H$ is tricyclic its credit is $2+D(H)/2>2$, which pays
$R$.  If their packing number is three, the triangles are pairwise
vertex-disjoint.  The reduced block tree joining them contains an actual
bridge; cut one such bridge so that the two retained induced cyclic territories
contain two triangles and one triangle.  Give the whole connector component at
the cut to one endpoint territory.  Their credits are $>1$ and $>0$, so
together with $R$ the total is strictly greater than $1+0-1=0$.  This exhausts
$k=0,1,2$ and specifies ownership of every connector, cut, and rooted branch.

For $\Theta(1,2,2)+C_3+C_5$, the shared-cut incidence tree
has either separate clusters or one central block.  Separating bridges use
the credits $>1$, $>0$, and $\geq2-\sqrt5$.  In one cluster, opening the
theta internal vertex away from its triangle cut leaves either a two-triangle
packet against a pentagonal packet, or a shared-triangle
$C_3,C_3,C_5$ packet against one tree.  Their credits are respectively
$>1+(2-\sqrt5)>0$ and $>(4-\sqrt5)-1>0$.  If the external triangle or
pentagon is the central block, split that cycle at one cut; the remaining path
belongs to the packet at the other cut.  These are all central-node choices
and assign every cut and branch once \cite{ThetaTwoCycles}.

For $2+2$, \eqref{eq:theta-class} shows that only two diamonds
$\Theta(1,2,2)$ can exceed the budget.  If an actual bridge separates them,
cut it and assign the complete connector to either endpoint territory; the two
attached-diamond credits are each $>1$.  Suppose instead that they share a cut
$w$, and keep $w$, the first diamond, and all branches assigned to that side
in a territory $H$.  If $w$ has degree three in the second diamond, its other
three vertices induce a path, which with their rooted branches is one nonempty
tree $T$.  If $w$ has degree two there, choose a degree-three neighbor $p$ and
transfer $p$, the edge $wp$, and every branch rooted at $p$ to $H$.  This adds
only a rooted tree at $w$, so $H$ is still an attached-diamond packet; the two
vertices left from the second diamond induce an edge and, with their branches,
form one nonempty tree $T$.  In both cut orbits the territories are disjoint,
connected, and induced, every cut and branch is owned exactly once, and
$\sig(H)+\sig(T)>1-1=0$.  This is the complete two-diamond argument
\cite{TwoDiamonds}.

For $3+1$, use the complete rank-three ledgers from the tricyclic proof
\cite{TricyclicGraphs}.  The following table makes the canonical/noncanonical
split exhaustive.  Here $j$ is the number of long paths in the all-odd
$K_4$ class; ``canonical pair'' means lengths $\{1,2\}$ in each doubled
pair.  Every number in the fourth column is an exact rational upper bound,
not a decimal optimization.
\begin{center}
\small
\begin{tabular}{p{2.4cm}p{3.3cm}p{3.0cm}p{3.2cm}}
\toprule
rank-three family&physical state&number/disposition&actual excess certificate\\
\midrule
four-path theta&all simple rows&all DNN&$e_B\leq2$; equality only at $(1,2,2,2)$\\
doubled triangle&outside class $111$&$28$ labelled rows, DNN&$e_B\leq2$ from the physical Gram ledger\\
doubled triangle, class $111$&a parallel path noncanonical&four relabelled choices, DNN&even-long $e_B<229/120<2$; odd-long $e_B<31/20<2$\\
doubled triangle, class $111$&both doubled pairs canonical&four labelled rows, structural&hostile-cycle certificate $e_B<221/100$\\
doubled $C_4$&outside class $111$&$28$ labelled rows, DNN&$e_B\leq2$ from the physical Gram ledger\\
doubled $C_4$, class $111$&a doubled path noncanonical&either pair and either parity, DNN&even-long $e_B<1862/1000<2$; odd-long $e_B<1662/1000<2$\\
doubled $C_4$, class $111$&both doubled pairs canonical&eight labelled rows, structural&hostile-cycle certificate $e_B<9/4$\\
$K_4$, non-all-odd&all physical rows&$56$ labelled rows, DNN&$e_B\leq5/3<2$\\
$K_4$, all odd&$j\geq3$&all subsets, DNN&simplex: $e_B<3-j/3\leq2$\\
$K_4$, all odd&$j=2$, opposite&three placements, DNN&$e_B\leq24-16\sqrt2<2$\\
$K_4$, all odd&$j=2$, adjacent&twelve placements, DNN&$e_B<71/40<2$\\
$K_4$, all odd&$j=1$&six placements, structural&hostile-cycle certificate $e_B<12/5$\\
$K_4$, all odd&$j=0$&the unsubdivided $K_4$, structural&attached-$K_4$ Sachs packet\\
\bottomrule
\end{tabular}
\end{center}
The two doubled-triangle bounds come from the planar angle records
$(12;16,7)$ and $(6;2,8)$ at physical lengths $(4,1,2,1,1)$ and
$(2,3,2,1,1)$.  The doubled-$C_4$ bounds come from
$Q(0,3,2,5)$ and $Q(0,1,1,4)$ at $(4,1,2,2,1,1)$ and
$(2,3,2,2,1,1)$.  Reflection and interchange within a doubled pair cover
every labelled placement, and fixed-parity monotonicity covers every further
length increase.  For all-odd $K_4$, the two-long paths are either adjacent or
opposite, and the $j\geq3$ simplex estimate covers every remaining subset.
Thus the $j=0,1$ rows, and only those all-odd rows, require structural
ownership.  This also verifies directly that an external triangle closes by
DNN every newly listed noncanonical state: each of
\[
 \frac{229}{120},\quad\frac{31}{20},\quad\frac{1862}{1000},
 \quad\frac{1662}{1000},\quad\frac{71}{40},\quad
 24-16\sqrt2
\]
is strictly below two, and $\epsilon_3=1$.

For a hostile cycle $Q=C_q$, where $q\equiv1\pmod4$ and $q\geq5$,
\[
 \epsilon_q\leq5-2\sqrt5<\frac35.
\]
All DNN rows in the table have $e_B\leq2$, so $e_B+\epsilon_q<3$.
For the three non-$K_4$ structural types and the one-long all-odd $K_4$ type,
the displayed bounds give respectively
$221/100+3/5<3$, $9/4+3/5<3$, and $12/5+3/5=3$ with strictness in both
constituent estimates.  The attached $K_4$ packet has credit $>2$; an actual
bridge separating it from $Q$ is cut and its whole connector is assigned to
one endpoint territory, giving total credit
$>2-(\sec(\pi/q)-1)>0$.  At a shared cut $z$, the $K_4$ territory owns $z$
and every branch rooted on its side, while the other territory owns $Q-z$ and
its branches; the latter is one nonempty tree, so the total is $>2-1>0$.

It remains to specify ownership for the structural rows when $Q$ is a triangle
or has length $3\pmod4$.  For a canonical doubled-triangle or doubled-$C_4$
row, open exactly the internal path vertex used in its tricyclic repair; for a
one-long all-odd $K_4$ row, open an internal vertex of the long path.  Give the
opened vertex every component whose unique block-cut route first meets the
rank-three block there.  The complementary territory contains the same
attached two-triangle or diamond packet as in the tricyclic proof.  If the
opened territory owns $Q$, it is a favorable unicyclic packet of credit $>0$
and the complement has credit $>1$.  Otherwise it is one nonempty tree of
credit $-1$, while the complement has three favorable cycles.  Packing number
at most two gives complement credit $>2$ by the Sachs argument above; packing
number three supplies an actual bridge separating credits $>1$ and $>0$.
Every connector remnant stays whole on one side, every cut has one owner, and
each rooted branch follows the territory owning its root.  For $j=0$, the
attached-$K_4$ territory owns the complete $K_4$ and its rooted branches; a
separating route leaves a favorable unicyclic territory, while at a shared cut
$z$ the other territory is the tree $Q-z$.  Even cycles have zero DNN excess.
This exhausts the table and assigns every vertex, connector, cut, and branch
exactly once; the underlying Gram records are reproduced in
\cite{RankThreeCycle}.  This proves every multiblock alternative.
\end{proof}

\section{Classification of the rank-four kernels}

It remains to consider one positive-rank block.  Suppress every maximal path
whose internal vertices have block degree two, retaining parallel kernel
edges.  The original graph is simple, so at most one path in a parallel class
can have length one.

\begin{proposition}[Seventeen-kernel classification]\label{prop:kernels}
Let $K$ be a loopless multigraph with at least two vertices, no cut vertex,
minimum degree at least three, and cyclomatic rank four.  Then
$2\leq|V(K)|\leq6$, and $K$ is isomorphic to exactly one of the seventeen
multigraphs in Table~\ref{tab:kernels}.
\end{proposition}

\begin{table}[p]
\centering\small
\begin{tabular}{r r l l r}
\toprule
No.&$v$&degree multiset&upper-triangle multiplicities&simplicity cost\\
\midrule
1&2&$5,5$&$5$&4\\
2&3&$5,4,3$&$1,2,3$&3\\
3&3&$4,4,4$&$2,2,2$&3\\
4&4&$5,3,3,3$&$0,1,2,1,2,1$&2\\
5&4&$4,4,3,3$&$0,1,2,2,1,1$&2\\
6&4&$4,4,3,3$&$0,1,2,2,2,0$&3\\
7&4&$4,4,3,3$&$0,1,2,3,1,0$&3\\
8&4&$4,4,3,3$&$1,1,1,1,1,2$&1\\
9&5&$4,3,3,3,3$&$0,0,1,2,1,0,2,2,0,0$&3\\
10&5&$4,3,3,3,3$&$0,0,1,2,1,1,1,1,1,0$&1\\
11&5&$4,3,3,3,3$&$0,0,1,2,1,1,1,2,0,0$&2\\
12&5&$4,3,3,3,3$&$0,1,1,1,1,1,1,0,1,1$&0\\
13&6&$3,3,3,3,3,3$&$0,0,0,1,2,0,1,1,1,2,1,0,0,0,0$&2\\
14&6&$3,3,3,3,3,3$&$0,0,0,1,2,0,1,2,0,2,0,1,0,0,0$&3\\
15&6&$3,3,3,3,3,3$&$0,0,0,1,2,1,1,0,1,1,1,0,1,0,0$&1\\
16&6&$3,3,3,3,3,3$&$0,0,1,1,1,0,1,1,1,1,1,1,0,0,0$&0\\
17&6&$3,3,3,3,3,3$&$0,0,1,1,1,1,0,1,1,1,0,1,1,0,0$&0\\
\bottomrule
\end{tabular}
\caption{Canonical kernel list.  Multiplicities use lexicographic
upper-triangle order; the final column is
$\sum_{u<v}\max(m_{uv}-1,0)$.}\label{tab:kernels}
\end{table}

\begin{proof}
Let $v=|V(K)|$ and $e=|E(K)|$.  Since $e=v+3$,
\begin{equation}\label{eq:degree-excess}
 \sum_x(d(x)-2)=2e-2v=6.
\end{equation}
Every summand is positive, so $2\leq v\leq6$.  The positive integers
$d(x)-2$ partition six.  For each resulting degree multiset introduce a
nonnegative integer $a_{ij}$ for every pair $i<j$ and solve
\[
 \sum_{j\ne i}a_{\min(i,j),\max(i,j)}=d_i.
\]
Reject a solution if its positive support disconnects after deletion of some
vertex, and quotient by permutations preserving the degree multiset.  The
possible degree rows and surviving orbit counts are
\[
\begin{array}{c|c|c}
v&\text{degree multisets}&\text{orbits}\\ \hline
2&(5,5)&1\\
3&(5,4,3),(4,4,4)&1+1\\
4&(5,3,3,3),(4,4,3,3)&1+4\\
5&(4,3,3,3,3)&4\\
6&(3,3,3,3,3,3)&5.
\end{array}
\]
Solving these incidence equations gives exactly the vectors in
Table~\ref{tab:kernels}.  Their lexicographically least vectors over all
labellings are distinct, proving nonisomorphism.  Suppression cannot create a
loop for the same cut-vertex reason used in Proposition~\ref{prop:block-alternatives}.
The independently regenerated list has canonical checksum
\begin{center}\ttfamily\small
d89e6e60c66e480ba89e662ab90b5ace211cbcff7292f92ad1614bb0937eb8e9.
\end{center}
\end{proof}

\section{Exact certificate ledgers for all seventeen kernels}

A physical row records, in each kernel bundle, the number of odd paths.  It is
not a switching class.  A certificate consists of an explicit correlation
Gram matrix on branch vertices and, where needed, exact rational planar
vectors on path interiors.  Its cost is the sum in \eqref{eq:kernel-min}.
Every accepted cost is at most the tetracyclic budget three.  A canonical row
uses length two for an even path and length one for an odd path, except that
parallel odd paths use lengths one and three to preserve simplicity.  A
one-coordinate frontier adds two to one physical path.  By
Lemma~\ref{lem:path}, the canonical certificate covers a row unchanged in all
coordinates, while any noncanonical descendant is covered by a frontier in
one changed coordinate and monotonicity in every other coordinate.

\begin{proposition}[Kernel certificate theorem]\label{prop:kernel-cert}
Every simple subdivision $H$ of every kernel in Table~\ref{tab:kernels}
satisfies
\[
 \kap(H)\leq|E(H)|+3
\]
or has an induced structural packet proving directly that
$\splus(H)\geq|V(H)|$.  The same conclusion holds after arbitrary rooted trees
are attached at arbitrary vertices.
\end{proposition}

\begin{proof}
We summarize the exact ledgers by branch-vertex count.  Each cited proof note
prints the formulas or defines the digest-locked fixture and a fail-closed
exact verifier.

\paragraph{Two branch vertices: kernel 1.}
The block is five internally disjoint terminal paths.  Its excess is the
five-term version of \eqref{eq:theta-excess}.  Fixed-parity monotonicity reduces
nonunit lengths to two or three.  If no path is unit, endpoint tests cover
four parity counts; in the remaining count the objective has negative
derivative at the endpoint value three.  If one path is unit, four parity
counts are handled by $\theta=\pi$ or $2\pi/3$.  In the last count, with
$u=\tan(\theta/4)$, the upper bound is
\[
 \frac1{4u^2}-\frac12+\frac{33}{4}u^2,
 \qquad \min=\frac{\sqrt{33}-1}{2}<3.
\]
Thus the excess is strictly below three for every simple five-path block
\cite{FivePath}.

\paragraph{Three branch vertices: kernels 2--3.}
Their bundle multiplicities are $(1,2,3)$ and $(2,2,2)$.  If $q$ of the six
paths are odd and $q\geq3$, the equicorrelation matrix with off-diagonal
$-1/2$ gives cost
\[
 \frac q3+\frac{2(6-q)}3=4-\frac q3\leq3.
\]
For $q\leq2$, the complete rational ledgers have respectively nine labelled
bundle-count rows and four rows up to bundle permutation.  All thirteen costs
are strict; their maxima are $21/8$ and $17/8$.  Every $3\times3$ Gram
determinant is listed as a nonnegative rational number \cite{ThreeVertex}.

\paragraph{Four branch vertices: kernels 4--8.}
The physical row counts are
\[
 72,\ 72,\ 54,\ 48,\ 96,
 \qquad\text{total }342.
\]
The exact disjoint ledger is
\begin{equation}\label{eq:four-ledger}
 342=270\ \text{base}+70\ \text{rational patch}
       +2\ \text{independently discharged}.
\end{equation}
The last two rows are kernel 8 with all five singleton bundles odd and the
doubled bundle containing one or two odd paths.  An exhaustive
$2\cdot2^6=128$ long/unit state census has 56 stabilizer orbits.  There are
114 states with at least two long support paths, twelve one-long states in six
planar Gram classes, and two no-long structural states.  A regular-simplex
matrix proves the first group strictly below three; exact planar angles with
denominators $3,5,7$ prove the second; in the last group, deleting the extra
path leaves an attached all-odd $K_4$ packet against one tree.  The verifier
proves all trigonometric comparisons with rational intervals and rejects eight
hostile mutations.  The canonical payload digest is
\begin{center}\ttfamily\small
f381b2b28bd3f45d7c96d90bce824a308bc340c9d6ebd098e5da116b89648d5a.
\end{center}
See \cite{FourVertex}.

\paragraph{Five branch vertices: kernels 9--12.}
The physical and genuine automorphism-orbit ledgers are
\[
\begin{array}{c|rrrr|r}
\text{kernel}&9&10&11&12&\text{total}\\ \hline
\text{physical rows}&108&192&144&256&700\\
\text{orbits}&63&120&144&51&378.
\end{array}
\]
An exact equilateral three-color sieve, checking all $3^5=243$ labelled
colorings per orbit, gives
\begin{equation}\label{eq:five-ledger}
 378=370\ \text{direct}+8\ \text{residual}.
\end{equation}
The physical odd--odd doubled-bundle cost uses lengths $(1,3)$ and the
algebraic number $a=3\tan^2(\pi/18)$, isolated as the unique root in
$(93/1000,94/1000)$ of
$a^3-27a^2+99a-9$.  Thus no floating comparison occurs.

The eight residual orbits split as $4,2,2,0$ across kernels 9--12.  Seven
non-symbolic rows use sixty exact rational records: canonical and
one-coordinate $+2$ frontiers for six eight-path rows, plus six doubled-path
frontiers for one kernel-9 row.  The remaining kernel-9 equality row uses five
planar vectors over $\mathbb Q(\sqrt3)$: two odd singleton paths have
correlation $-1$ and cost zero, while each of three mixed doubled bundles has
correlation $-1/2$ and cost $1/3+2/3=1$.  Its total is exactly three.
The closure verifier recomputes every rational internal-vector ledger and the
symbolic Gram identities and rejects ten hostile mutations
\cite{FiveSieve,FiveClosure}.

\paragraph{Six branch vertices: kernels 13--17.}
Kernel 16 is $K_{3,3}$.  Its $2^9=512$ physical parity rows have exact
three-color score histogram
\[
\begin{array}{c|rrrrrrrr}
S&0&3&4&5&6&7&8&9\\ \hline
\#&1&6&27&63&147&168&81&19.
\end{array}
\]
If $q$ odd paths and $b$ even paths are bichromatic, the excess is
$(q+2b)/3=S/3\leq3$.  The verifier checks all $3^6=729$ colorings for every
row and all 26 genuine automorphism orbits \cite{KernelSixteen}.

For kernels 13, 14, 15, and 17, the labelled physical-row counts are
$288,216,384,512$, the automorphism-group orders are $4,6,4,12$, and the
genuine orbit counts are $102,56,144,74$.  Thus there are $1400$ labelled
rows and 376 genuine orbits.  The exact partition is
\begin{equation}\label{eq:six-sieve}
 376=359\ \text{three-color certified}+17\ \text{residual},
 \qquad (5,6,5,1)\text{ residual by kernel}.
\end{equation}
For the sixteen residual orbits of kernels 13--15, canonical and nine
one-coordinate frontiers give
\begin{equation}\label{eq:six-frontier}
 160=148\ \text{strict rational certificates}
        +12\ \text{symbolic kernel-14 targets}.
\end{equation}
The twelve targets are the canonical row and coordinates $0,3,8$ for each of
three kernel-14 rows.  Their three Gram determinants are $1/2,0,1/2$; every
$2\times2$ principal minor is $3/4$, each mixed doubled bundle costs one, and
each canonical total is exactly three.

Kernel 17 has 512 labelled parity rows and 74 genuine automorphism orbits.
Seven exact planar templates over $\mathbb Q(\sqrt3)$, with all 32 branch sign
switches per template, cover all 512 physical rows.  Their disjoint first-cover
gains are
\[
 284,\ 123,\ 56,\ 24,\ 13,\ 8,\ 4,
\]
which sum to 512.  In particular, they cover the residual orbit at its
canonical vector and all nine coordinate frontiers.  The final verifier
requires the 359-row sieve, 148-record fixture, twelve symbolic targets, and
seven-template audit, then reconstructs all ledgers.  The residual is zero
\cite{CubicFinal}.

All direct rows now satisfy \eqref{eq:budget}.  Every structural row uses an
induced partition whose deleted territory and all branches based there are
owned explicitly; retained branches stay with the complementary packet.
Fixed-parity monotonicity proves arbitrary subdivisions, and
Lemma~\ref{lem:dnn} proves arbitrary rooted-tree attachments.  This completes
all seventeen kernels.
\end{proof}

\begin{proof}[Proof of Theorem~\ref{thm:main}]
Apply Proposition~\ref{prop:block-alternatives}.  If there is more than one
positive-rank cyclic block, Proposition~\ref{prop:multiblock} applies.  If
there is one rank-four block, suppress its degree-two paths.  Proposition
\ref{prop:kernels} gives one of the seventeen kernels, and Proposition
\ref{prop:kernel-cert} proves the required inequality for every subdivision
and every attached rooted tree.  The alternatives are exhaustive, so
$\splus(G)\geq|V(G)|$.
\end{proof}

\appendix
\section{Verifier manifest and exact scopes}\label{app:verify}

All commands below are run from the repository root.  Run every finite audit
normally and with Python assertions disabled; the paired outputs must be
byte-identical where the verifier says so.

\begin{verbatim}
python3 research/rank-four-kernel-census-verifier.py
python3 -O research/rank-four-kernel-census-verifier.py

python3 research/rank-four-three-vertex-tables-verifier.py
python3 -O research/rank-four-three-vertex-tables-verifier.py

python3 research/rank-four-four-vertex-theorem-verifier.py
python3 -O research/rank-four-four-vertex-theorem-verifier.py

python3 research/rank-four-five-vertex-three-color-verifier.py
python3 -O research/rank-four-five-vertex-three-color-verifier.py
python3 research/rank-four-five-vertex-residual-closure-verifier.py
python3 -O research/rank-four-five-vertex-residual-closure-verifier.py

python3 research/rank-four-kernel16-three-color-verifier.py
python3 -O research/rank-four-kernel16-three-color-verifier.py
python3 research/rank-four-cubic-kernels-final-verifier.py
python3 -O research/rank-four-cubic-kernels-final-verifier.py

python3 research/tetracyclic-master-verifier.py
python3 -O research/tetracyclic-master-verifier.py
\end{verbatim}

The census verifier independently enumerates weak edge-multiplicity
compositions on two through six vertices.  It finds respectively
$1,7,54,255,550$ labelled survivors and collapses them to
$1,2,5,4,5$ canonical classes.  It checks the exact list, degree excess,
simplicity costs, checksum, and hostile mutations.

The three-vertex verifier reconstructs all $24+27=51$ physical rows, split as
$9+10$ low-odd rational-table rows and $15+17$ common-simplex rows, and checks
every Gram principal minor exactly.  The four-vertex verifier reconstructs the disjoint
$342=270+70+2$ ledger and all 128 long/unit frontier states.  The five-vertex
pair reconstructs all 700 physical rows, 378 orbits, the
$370+8$ sieve, the sixty rational records, and the symbolic equality row.  The
kernel-16 verifier reconstructs all 512 rows and its score histogram.  The
final cubic verifier invokes its dependency audits and checks
$376=359+17$, $160=148+12$, all symbolic kernel-14 matrices, and all 512
kernel-17 first-cover assignments.  The master verifier invokes all seven
finite dependency verifiers, verifies their file and output digests, checks
that kernels $1$ through $17$ have exactly one theorem owner (with the stated
sieve/closure composition for kernels $9$--$12$), and records kernel 1 as the
separate analytic five-path theorem.

These programs are exact audits of finite combinatorial and rational or
quadratic-algebraic proof objects.  They do not replace the mathematical proofs
of DNN duality, path elimination, monotonicity, block classification, induced
superadditivity, or the multiblock territory lemmas.  Conversely, the paper
does not ask a floating-point optimizer or finite graph-order census to prove
an all-length assertion.

The detailed proof-object sources are:
\begin{verbatim}
positive-square-energy/tetracyclic-general/
  five-path-dnn-theorem.md
  three-vertex-kernels-dnn.md
  four-vertex-kernels-dnn.md
  five-vertex-three-color-dnn-sieve.md
  five-vertex-residual-closure-theorem.md
  kernel16-three-color-dnn.md
  cubic-kernels-13-15-17-final-theorem.md
  rank-four-kernel-classification.md
  theta-two-cycle-residuals.md
  two-diamond-blocks.md
  rank3-hostile-cycle-packet.md
\end{verbatim}

Build this manuscript with
\begin{verbatim}
bash scripts/build-paper.sh all-tetracyclic-graphs
\end{verbatim}

\paragraph{Scoped nonclaim.}
The theorem concerns exactly finite simple connected graphs with $m=n+3$.
The argument does not assume that adding an edge increases $\splus$, does not
transport physical path lengths through switching, and does not infer a
quantitative reserve from an unquantified lower-rank theorem.  Several DNN
certificates attain their auxiliary excess budget three.  We therefore assert
only $\splus(G)\geq n$ globally, not strict inequality and not a classification
of spectral equality.

\section*{AI disclosure}
This manuscript and its proof-object synthesis were generated by an AI coding
and mathematical reasoning system from the cited repository materials.  No
human authorship, review, or verification is claimed.  The result should be
assessed from the displayed argument and exact local dependencies.  No
external publication was performed as part of this work.

\begin{thebibliography}{99}
\bibitem{BicyclicGraphs}
Companion manuscript, \emph{Positive Square Energy of Every Bicyclic Graph},
2026; \path{all-bicyclic-graphs/paper.tex}.

\bibitem{TricyclicGraphs}
Companion manuscript, \emph{Positive Square Energy of Every Tricyclic Graph},
2026; \path{all-tricyclic-graphs/paper.tex}.

\bibitem{TetracyclicCactus}
Companion manuscript, \emph{Positive Square Energy of Every Tetracyclic
Cactus}, 2026; \path{all-tetracyclic-cacti/paper.tex}.

\bibitem{FivePath}
\emph{The five-path rank-four block}, proof note, 2026;
\path{positive-square-energy/tetracyclic-general/five-path-dnn-theorem.md}.

\bibitem{ThreeVertex}
\emph{Three-vertex rank-four kernels: a uniform DNN certificate}, proof note,
2026; \path{positive-square-energy/tetracyclic-general/three-vertex-kernels-dnn.md}.

\bibitem{FourVertex}
\emph{Four-vertex rank-four kernels: complete exact ledger}, proof note, 2026;
\path{positive-square-energy/tetracyclic-general/four-vertex-kernels-dnn.md}.

\bibitem{FiveSieve}
\emph{Five-vertex rank-four kernels: exact coarse three-color sieve}, proof
note, 2026;
\path{positive-square-energy/tetracyclic-general/five-vertex-three-color-dnn-sieve.md}.

\bibitem{FiveClosure}
\emph{Five-vertex rank-four kernels: closure of the eight residual rows}, proof
note, 2026;
\path{positive-square-energy/tetracyclic-general/five-vertex-residual-closure-theorem.md}.

\bibitem{KernelSixteen}
\emph{Kernel 16: exact three-color certificate on all physical rows}, proof
note, 2026;
\path{positive-square-energy/tetracyclic-general/kernel16-three-color-dnn.md}.

\bibitem{CubicFinal}
\emph{Cubic kernels 13--15 and 17: exact final theorem}, proof note, 2026;
\path{positive-square-energy/tetracyclic-general/cubic-kernels-13-15-17-final-theorem.md}.

\bibitem{ThetaTwoCycles}
\emph{The tetracyclic $2+1+1$ residuals}, proof note, 2026;
\path{positive-square-energy/tetracyclic-general/theta-two-cycle-residuals.md}.

\bibitem{RankThreeCycle}
\emph{Rank-three block plus one hostile cycle: canonical residual packet},
proof note, 2026;
\path{positive-square-energy/tetracyclic-general/rank3-hostile-cycle-packet.md}.

\bibitem{TwoDiamonds}
\emph{Two diamond blocks}, proof note, 2026;
\path{positive-square-energy/tetracyclic-general/two-diamond-blocks.md}.
\end{thebibliography}

\end{document}
